% Authority: German Wikiversity pageid 151178, revid 1095920, 2026-06-15T07:26:34Z.
% Exact UTF-8 wikitext witness SHA-256: 237b3b08df28c060d3c5c2fb6efdf22ea7ae5df6c1fbfad68d64380fe7b4f852.
% Expanded German TeX source SHA-256: ed49f1889840b8352f634ef01400301849f8303ec27a2a38b334b744ae5e5951.

\definitionsverweis {Medan gradien}{}{}
dari $h$ adalah $\begin{pmatrix}  2x \\ 2y   \end{pmatrix}$, sehingga sebuah medan normal satuan adalah

\mavergleichskettedisp
{\vergleichskette
{ N \begin{pmatrix}  x  \\ y   \end{pmatrix}
}
{ =} { { \frac{   1 }{  \sqrt{x^2+y^2}  } } \begin{pmatrix}  x  \\ y   \end{pmatrix}
}
{ } {
}
{ } {
}
{ } {
}
}
{}{}{.}
\definitionsverweis {Diferensial total}{}{}
$N$ di suatu titik

\mavergleichskettedisp
{\vergleichskette
{ P
}
{ =} { \begin{pmatrix}  x  \\ y   \end{pmatrix}
}
{ } {
}
{ } {
}
{ } {
}
}
{}{}{}
adalah

\mavergleichskettedisp
{\vergleichskette
{  \left(D N \right)_{P}
}
{ =} { \begin{pmatrix}  { \frac{   y^2 }{   \left(  x^2+y^2  \right)^{3/2}  } }   &   - { \frac{   xy }{   \left(  x^2+y^2  \right)^{3/2}  } }   \\    - { \frac{   xy }{   \left(  x^2+y^2  \right)^{3/2}  } }    &  { \frac{   x^2 }{   \left(  x^2+y^2  \right)^{3/2}  } }  \end{pmatrix}
}
{ } {
}
{ } {
}
{ } {
}
}
{}{}{.}
Penerapan pada vektor singgung
\mathl{\begin{pmatrix}  y  \\ -x   \end{pmatrix}}{}
memberikan

\mavergleichskettealign
{\vergleichskettealign
{ \left(D N \right)_{P} \begin{pmatrix}  y  \\ -x   \end{pmatrix}
}
{ =} { \begin{pmatrix}  y^2    &   - xy    \\  - xy   &  x^2  \end{pmatrix} \begin{pmatrix}  y  \\ -x   \end{pmatrix}
}
{ =} { \begin{pmatrix}  y \left(  y^2 +x^2  \right)  \\ -x \left(  y^2 +x^2  \right)    \end{pmatrix}
}
{ =} { \begin{pmatrix}  y   \\ -x    \end{pmatrix}
}
{ } {
}
}
{}
{}{.}
Oleh karena itu, menurut
[[Ebene differenzierbare Kurve/Faser/Krümmung/Fakt|fakta]]
kelengkungan kurva di $P$ sama dengan $-1$.
\zusatzklammer {Hal ini sesuai karena
\mathl{\begin{pmatrix}  y  \\ -x   \end{pmatrix}}{} dan medan gradien merepresentasikan orientasi standar, sedangkan
\mathl{\begin{pmatrix}  y  \\ -x   \end{pmatrix}}{} merupakan arah singgung untuk gerak searah jarum jam.} {} {}

 [[Kategorie:Latexseite]]
